\documentclass[12pt]{article}
\usepackage{sbc-template}
\usepackage[utf8]{inputenc}
\usepackage[T1]{fontenc}
\usepackage[brazil]{babel}
\usepackage{graphicx}
\usepackage{url}

\title{Trabalho Prático I - AEDS III:\\Manipulação de Arquivo Sequencial e Ordenação Externa}

\author{Bruno Pais Moacir}

\address{Instituto de Ciências Exatas e Informática (ICEI) \\
  Pontifícia Universidade Católica de Minas Gerais (PUC Minas)\\
  Belo Horizonte -- MG -- Brasil
  \email{bruno@sga.pucminas.br} % Ajuste para seu e-mail acadêmico real
}

\begin{document}

\maketitle

\begin{abstract}
This report details the implementation of a sequential file management system developed for the Data Structures and Algorithms III (AEDS III) course. The project involves parsing a public domain dataset, performing CRUD operations directly on a binary file using tombstone logic for logical deletion, and executing an external sorting algorithm (Balanced Multiway Merging). The core constraints included avoiding high-level native string manipulation methods, ensuring structural robustness, and managing variable-length records in secondary memory.
\end{abstract}

\begin{resumo}
Este relatório detalha a implementação de um sistema de gerenciamento de arquivos sequenciais desenvolvido para a disciplina de Algoritmos e Estruturas de Dados III (AEDS III). O projeto envolve a leitura de uma base de dados de domínio público, a execução de operações CRUD diretamente em um arquivo binário utilizando a lógica de lápides para exclusão lógica, e a execução de um algoritmo de ordenação externa (Intercalação Balanceada). As principais restrições incluíram evitar métodos nativos de alto nível para manipulação de strings, garantir a robustez estrutural e gerenciar registros de tamanho variável em memória secundária.
\end{resumo}

\section{Introdução}
O Trabalho Prático I de Algoritmos e Estruturas de Dados III tem como objetivo consolidar os conceitos de armazenamento de dados em memória secundária. A manipulação de arquivos sequenciais requer uma gestão rigorosa do espaço, especialmente quando os registros possuem tamanhos variados. Para este projeto, foi selecionada a base de dados \textit{IMDb Movies}, contemplando campos obrigatórios como strings de tamanho variável (título do filme), strings de tamanho fixo (sigla do país), datas de lançamento, listas de valores (gêneros) e valores de ponto flutuante (avaliações).

\section{Desenvolvimento}
O sistema foi desenvolvido integralmente em Java, com forte ênfase na clareza de codificação e obediência às restrições arquiteturais propostas.

\subsection{Carga e Parse Manual de Dados}
A importação da base de dados CSV para o arquivo binário foi construída de forma estrita. Para evitar violações das regras da disciplina, bibliotecas ou métodos nativos complexos como \texttt{.split()} e \texttt{.replaceAll()} foram descartados. A separação de colunas, a extração da lista de gêneros e a conversão de datas foram programadas manualmente utilizando laços de repetição e o método \texttt{.charAt()}, garantindo total controle sobre a varredura das strings.

\subsection{Estrutura do Arquivo Binário e CRUD}
A memória secundária (\texttt{dados.bin}) foi estruturada com um cabeçalho inicial de 4 bytes contendo o último ID gerado. Cada registro gravado sequencialmente obedece à estrutura: Lápide (1 byte), Indicador de Tamanho (4 bytes) e o Vetor de Bytes do objeto. 
As operações de CRUD lidam com a movimentação do ponteiro via \texttt{RandomAccessFile}. Destaca-se a operação de Atualização (\textit{Update}): caso o novo vetor de bytes exceda o tamanho do registro original, o registro antigo recebe a marcação de exclusão lógica (lápide com \texttt{*}) e o registro atualizado é movido para o fim do arquivo.

\subsection{Ordenação Externa}
Para lidar com a desordenação causada pelos \textit{updates} e para recuperar o espaço dos registros deletados logicamente, foi implementado o algoritmo de Intercalação Balanceada. O processo foi dividido em duas fases:
\begin{itemize}
    \item \textbf{Distribuição:} O arquivo original é lido em blocos (limitados pela capacidade da memória primária). Os registros válidos são ordenados por ID e gravados em arquivos temporários (caminhos).
    \item \textbf{Intercalação:} Os arquivos temporários são lidos simultaneamente. O menor registro entre os caminhos é selecionado e gravado no arquivo final definitivo, garantindo a ordenação completa e a limpeza dos fragmentos deletados.
\end{itemize}

\section{Testes e Resultados}
Os testes foram conduzidos via terminal de linha de comando no ambiente Linux. A carga inicial gravou mais de 10.000 registros com sucesso. Operações de leitura retornaram os objetos formatados instantaneamente. Testes críticos de atualização provaram que o sistema gerencia corretamente o crescimento do vetor de bytes, validado pela busca imediata do registro após a exclusão do seu espaço original. Por fim, a Ordenação Externa foi executada com sucesso utilizando 3 caminhos e um limite de 20.000 registros, gerando um arquivo limpo e perfeitamente ordenado, sem corromper a integridade dos dados.

\section{Conclusão}
A implementação atingiu 100\% de conformidade com as especificações. A construção manual do \textit{parser} de strings e a estruturação em baixo nível do arquivo binário garantiram a robustez do programa. A correta execução da ordenação externa comprovou a eficiência da arquitetura escolhida para a gestão de memória secundária, preparando a base para as futuras implementações de indexação e árvores B+ nas próximas etapas da disciplina.

\end{document}